% Working draft for ACL 2027 (ACL Rolling Review, January 2027 cycle).
% Long paper: up to 8 pages of content in review mode; Limitations section required (not counted).
\documentclass[11pt]{article}
\usepackage[review]{acl}
\usepackage{times}
\usepackage{latexsym}
\usepackage[T1]{fontenc}
\usepackage[utf8]{inputenc}
\usepackage{microtype}
\usepackage{inconsolata}
\usepackage{graphicx}
\usepackage{booktabs}
\usepackage{amsmath}
\usepackage{xcolor}

% Numbers that depend on experiments still running (Phase A) are marked, never filled in.
\newcommand{\pending}[1]{\textcolor{red}{[#1]}}

\title{Allocating Recurrent Depth in Looped Language Models for Multi-Step Tool Use}

\author{Anonymous ACL submission}

\begin{document}
\maketitle

\begin{abstract}
Looped language models apply a shared stack of layers several times per token, so a single
model can trade inference compute for accuracy by changing how many loops it runs. This makes
them a natural fit for agents, whose requests range from single tool calls to long chains of
dependent calls. In practice, however, deployed looped models run a fixed number of loops for
every request, and per-token exit gates lose most of their savings once requests are decoded in
batches. Can a looped model decide, request by request, when extra loops are worth their cost?
We study depth allocation in looped models for multi-step tool use and propose an escalation
router: every request starts at a low loop count, and the model switches to full depth only when
its own early loops disagree on the decision tokens of its first tool call, a signal it computes
anyway. A controlled analysis motivates this design: choosing depth separately for each tool call
gives no gain beyond a sampling control with the same budget, whereas the depth a request needs
is a stable property that early loops reveal and the request text does not. On held-out NESTful
requests with a fine-tuned Ouro-2.6B, the router \pending{matches full-depth accuracy at X\% of
its compute}. We also show that most MathQA validation problems have near-duplicates in its
training split, a source of inflated results for work that trains on MathQA, and provide clean
splits.
\end{abstract}

\section{Introduction}
\label{sec:intro}

A looped (or recurrent-depth) language model applies a shared block of layers several times per
token instead of stacking distinct layers
\citep{dehghani2019universal,geiping2025recurrent,zhu2025ouro}. The loop count becomes an
inference-time dial: the same checkpoint can run shallow and cheap or deep and more accurate.
Looped models are now released by several groups, including Ouro \citep{zhu2025ouro}, Dory
\citep{nvidia2026dory} and Nanbeige4.2 \citep{nanbeige2026}, and looping has been shown to help
compositional tool calling specifically \citep{popescu2026toolcalling}. How a deployed model
should set the dial is much less settled.

Two practices dominate. Most deployments fix one loop count for all traffic; Dory's authors, for
example, recommend two loops as a default and four when accuracy matters more
\citep{nvidia2026dory}. The alternative is per-token halting, in the line of adaptive computation
time \citep{graves2016adaptive}, Universal Transformers \citep{dehghani2019universal} and
PonderNet \citep{banino2021pondernet}: Ouro trains an exit gate that can stop each token early
\citep{zhu2025ouro}, and recurrent-depth models can stop when successive loops stop changing the
prediction \citep{geiping2025recurrent}. Per-token exits fit poorly with how models are served.
Requests are decoded in batches, and a batch step costs as much as its deepest token; in our
measurements Ouro's gate is faster than a fixed depth of matching accuracy at batch size 1 but
slower from batch size 4 on, while a depth fixed per request keeps its savings under batching (\S\pending{serving}). Recent
analyses also find that learned halting gates do not track problem difficulty well
\citep{popescu2026halting,olszowiec2026monitoring}.

This paper asks at what granularity, and from what signal, a looped model should choose its
depth, using multi-step tool use as the test bed. Tool use suits the question: a request unfolds
as a sequence of calls, each call is executed, and the final result can be checked exactly, so
every allocation policy can be scored by accuracy and by the compute it spends. We fine-tune
Ouro-2.6B for chained calls, evaluate on the math portion of NESTful \citep{basu2024nestful}
with every intermediate call executed, and compare three granularities of allocation (per tool
call, per request and per token) against fixed depths and against random mixes of fixed depths
at equal compute. Confirmatory tests are fixed in advance and run once on held-out NESTful
requests.

Our findings are:

\begin{itemize}
  \item \textbf{Per-call depth is search, not allocation.} An oracle that picks the best loop
  count for each call reaches .308 on NESTful against .225 for four loops (with an earlier
fine-tuned checkpoint), which looks like a
  large gain. But an oracle that spends the same budget on sampled continuations at four loops
  reaches .325. The same holds for Dory-1.0-4B (.633 vs .617). Varying depth per call mostly
  perturbs the trajectory; it does not allocate compute where it is needed.
  \item \textbf{The useful decision is two loops or four.} On 240 NESTful requests, one loop
  solves .171, two .383, three .442 and four .492 of requests. Three loops is exactly as accurate
  as randomly mixing two- and four-loop runs at the same compute (.441), and one loop is too
  shallow even for single calls. Per request, 85 are solved at both two and four loops, 33 only
  at four and 7 only at two; a router that knew these outcomes would reach .521 at 0.83 times the
  compute of four loops.
  \item \textbf{Depth need is predictable from early loops, not from the request.} Whether a
  request is solved at a given depth is largely stable across sampled runs, yet features of the
  request text predict it no better than chance. Signals from the model itself do, once they are
  read at the right tokens: early loops agree with the final loop on the boilerplate that opens
  a tool call and disagree on the function name and argument values.
  \item \textbf{A deployable escalation router.} Every request starts at two loops; after the
  first ${\sim}45$ tokens, loop-1 and loop-2 signals on those decision tokens decide whether to
  switch to four loops. Nothing extra is computed to decide. \pending{Final held-out result:
  accuracy and compute relative to four loops on NESTful requests 240--719; gain over the best
  random mix with 90\% interval.}
  \item \textbf{A data-leakage audit.} Most MathQA validation problems (1,001 of 1,293 after our
  task construction) have a near-duplicate or an identical computation in the training split.
  Scores on them overstate held-out accuracy (at four loops, .489 on near-duplicates vs .346 on
  clean problems in one evaluation slice), which matters for any work that trains on MathQA and evaluates on
  NESTful, whose math portion comes from the MathQA test split.
\end{itemize}

\section{Related Work}
\label{sec:related}

\paragraph{Looped and recurrent-depth language models.}
Weight sharing across depth goes back to Universal Transformers \citep{dehghani2019universal}.
Recent work scales it to pretrained language models: a recurrent-depth model trained with a
random number of iterations can spend more loops at test time \citep{geiping2025recurrent};
Ouro trains a looped model with an entropy-regularised exit distribution over four loops
\citep{zhu2025ouro}; Dory-1.0 loops a middle block between fixed input and output blocks and
supports a last-loop KV cache \citep{nvidia2026dory}; Nanbeige4.2 runs its whole stack twice and
reports little benefit from more passes \citep{nanbeige2026}. Theory and controlled studies relate
loops to multi-step reasoning \citep{saunshi2025latent,kohli2026loop}, and fine-tuned looped
models improve compositional tool calling \citep{popescu2026toolcalling}. These works fix the
loop count or use the model's own halting rule at inference. Our work differs from
\citet{popescu2026toolcalling}, the closest to it, in question and in protocol: we study where
to spend loops rather than whether they help; calls are executed one turn at a time rather than
generated as one sequence; every allocation policy is compared with random mixes of fixed depths
and with a sampling control at equal compute, under batched serving; and confirmatory results
are fixed in advance and measured on held-out requests.

\paragraph{Adaptive depth: halting, early exit and routing.}
Adaptive computation time \citep{graves2016adaptive} and PonderNet \citep{banino2021pondernet}
learn when a recurrent computation should stop; Ouro's exit gate follows this line
\citep{zhu2025ouro}. Diagnoses of learned gates in looped transformers find that they are not
Pareto-optimal and that their depth use does not track problem difficulty \citep{popescu2026halting,olszowiec2026monitoring}, and Ouro's per-loop KV caches remove most of the
savings of exiting early \citep{mishra2026dynamic}. In standard
transformers, early exit stops at an intermediate layer when the model is confident
\citep{schuster2022calm,elhoushi2024layerskip}, and learned routers skip or repeat layers per
token or per input
\citep{raposo2024mod,he2025routertuning,luo2025flexidepth,heakl2026drllm,li2025cola}, in one
case choosing exit layers with reinforcement learning \citep{zhu2026lede}.
Mixture-of-Recursions trains a per-token router over recursion depth during pretraining
\citep{bae2025mor}. Our router differs in three ways: it acts once per request rather than per
token, so batched serving keeps its savings; it is fitted after training on a frozen looped model;
and it is evaluated on multi-step tool use with executed calls.

\paragraph{Speculative decoding with loops.}
Because a shallow loop's prediction is a prefix of the full computation, looped models can
draft tokens at few loops and verify them at full depth without changing the output
\citep{leviathan2023speculative,geiping2025recurrent}. LoopSpec pipelines this and reports about
3$\times$ faster batch-size-1 decoding on Ouro-2.6B \citep{cho2026loopspec}, with related
schemes in \citet{ha2026wavefront}, \citet{li2026das} and \citet{liu2026loopahead}. These methods
reduce latency at a fixed depth and leave the output unchanged; they do not decide the depth.
They are complementary to ours, and their observation that one-loop drafts usually match the
four-loop output agrees with our finding that loops disagree mainly on decision tokens.
Cross-loop distillation instead trains shallow loops to imitate the deepest one
\citep{wang2026loopopd}.

\paragraph{Allocating test-time compute.}
Spending more inference compute on harder inputs can beat a larger model at equal cost
\citep{snell2024scaling}, and learned predictors can allocate decoding budget per input
\citep{damani2024learning}. Width-based methods sample several answers and vote
\citep{wang2023selfconsistency}. Model routing and cascades choose between models of different
sizes from the input or from a cheap model's confidence
\citep{chen2024frugalgpt,ong2025routellm,gupta2024cascades}. We allocate depth within one looped
model, compare depth with width at equal compute, and show that the input alone does not reveal
which requests need more depth.

\paragraph{Tool use.}
Language models learn to call tools from demonstrations and synthetic data
\citep{schick2023toolformer,patil2024gorilla,liu2024apigen} and are evaluated on single and
parallel calls \citep{patil2025bfcl} and on nested sequences whose calls depend on earlier
results \citep{basu2024nestful}. NESTful's math portion is built from the MathQA test split
\citep{amini2019mathqa}, whose problems derive from AQuA-RAT \citep{ling2017aqua}; we use the
same construction on MathQA's training split for fine-tuning and audit its overlap with
validation and test.


% Sections still to write: setup (models, data, protocol, compute accounting, leakage audit),
% per-step allocation, what each loop count buys, predictability of depth need, the escalation
% router and its held-out confirmation, serving under batching, limitations.

\section*{Limitations}
\pending{To be written; must cover: single task family until Phase B, model sizes, greedy
decoding only, compute counted in loop-tokens until Phase D, outcome noise in per-request labels.}

\bibliography{references}

\end{document}
